\documentclass[10pt,a4paper]{amsart}
\usepackage{amsmath,amssymb,amsthm,graphicx}
% Retained minimal preamble for the Stacks Project chapter more-algebra.tex (More on Algebra).
% The chapter file pinned-original-source.tex carries only \input{preamble} on its first line, so no
% preamble text lives inside the chapter itself.  The Stacks Project uses one single shared preamble for
% every chapter; the snapshot stages it identically as sources/stacks-chow/preamble.tex and
% sources/stacks-etale-cohomology/preamble.tex.  The block below is copied mechanically, line for line,
% from sources/stacks-chow/preamble.tex lines 14-18 (the Xy-pic setup) and 154-217 (the native theorem
% declarations together with the mathematical macros).  Nothing was added, removed or redefined; the
% wrapper class and the amsmath/amssymb/amsthm/graphicx base packages are supplied by the assembler.
% Chapter: more-algebra.tex -- More on Algebra.
% For commutative diagrams we use Xy-pic
\usepackage[all]{xy}

% We use 2cell for 2-commutative diagrams.
\xyoption{2cell}

% Theorem environments.
%
\theoremstyle{plain}
\newtheorem{theorem}[subsection]{Theorem}
\newtheorem{proposition}[subsection]{Proposition}
\newtheorem{lemma}[subsection]{Lemma}

\theoremstyle{definition}
\newtheorem{definition}[subsection]{Definition}
\newtheorem{example}[subsection]{Example}
\newtheorem{exercise}[subsection]{Exercise}
\newtheorem{situation}[subsection]{Situation}

\theoremstyle{remark}
\newtheorem{remark}[subsection]{Remark}
\newtheorem{remarks}[subsection]{Remarks}

\numberwithin{equation}{subsection}

% Macros
%
\def\lim{\mathop{\mathrm{lim}}\nolimits}
\def\colim{\mathop{\mathrm{colim}}\nolimits}
\def\Spec{\mathop{\mathrm{Spec}}}
\def\Hom{\mathop{\mathrm{Hom}}\nolimits}
\def\Ext{\mathop{\mathrm{Ext}}\nolimits}
\def\SheafHom{\mathop{\mathcal{H}\!\mathit{om}}\nolimits}
\def\SheafExt{\mathop{\mathcal{E}\!\mathit{xt}}\nolimits}
\def\Sch{\mathit{Sch}}
\def\Mor{\mathop{\mathrm{Mor}}\nolimits}
\def\Ob{\mathop{\mathrm{Ob}}\nolimits}
\def\Sh{\mathop{\mathit{Sh}}\nolimits}
\def\NL{\mathop{N\!L}\nolimits}
\def\CH{\mathop{\mathrm{CH}}\nolimits}
\def\proetale{{pro\text{-}\acute{e}tale}}
\def\etale{{\acute{e}tale}}
\def\QCoh{\mathit{QCoh}}
\def\Ker{\mathop{\mathrm{Ker}}}
\def\Im{\mathop{\mathrm{Im}}}
\def\Coker{\mathop{\mathrm{Coker}}}
\def\Coim{\mathop{\mathrm{Coim}}}

% Boxtimes
%
\DeclareMathSymbol{\boxtimes}{\mathbin}{AMSa}{"02}

%
% Macros for moduli stacks/spaces
%
\def\QCohstack{\mathcal{QC}\!\mathit{oh}}
\def\Cohstack{\mathcal{C}\!\mathit{oh}}
\def\Spacesstack{\mathcal{S}\!\mathit{paces}}
\def\Quotfunctor{\mathrm{Quot}}
\def\Hilbfunctor{\mathrm{Hilb}}
\def\Curvesstack{\mathcal{C}\!\mathit{urves}}
\def\Polarizedstack{\mathcal{P}\!\mathit{olarized}}
\def\Complexesstack{\mathcal{C}\!\mathit{omplexes}}
% \Pic is the operator that assigns to X its picard group, usage \Pic(X)
% \Picardstack_{X/B} denotes the Picard stack of X over B
% \Picardfunctor_{X/B} denotes the Picard functor of X over B
\def\Pic{\mathop{\mathrm{Pic}}\nolimits}
\def\Picardstack{\mathcal{P}\!\mathit{ic}}
\def\Picardfunctor{\mathrm{Pic}}
\def\Deformationcategory{\mathcal{D}\!\mathit{ef}}
\title{The ramification index of a purely inseparable extension of discretely valued fields is a power of the characteristic}
\author{The Stacks Project Authors}
\date{stacks-project:more-algebra.tex:local-snapshot-20261012; GNU FDL 1.2 or later}
\begin{document}
\maketitle
\setcounter{subsection}{1}
\noindent\textit{Editorial note.} All delivered mathematical text is the author's own, copied from the pinned Stacks Project chapter file more-algebra.tex; no statement or proof step was rewritten, reordered or removed. The counted statement is the labeled lemma environment at source lines 32964-32970 (canonical tag 09E6) and the counted proof is the proof environment at source lines 32972-32978, which immediately follows it and is its own proof.\par\medskip
\noindent\textit{Editorial note.} The lemma immediately before this one, whose proof occupies source lines 32959-32962 and refers to a labelled item of another chapter, lies outside the counted spans; the discussion of unramified extensions that follows at source lines 32980-32983 is also outside them.\par\medskip
\noindent\textit{Editorial note.} The chapter file pinned-original-source.tex carries only an input of the shared preamble on its first line, so no preamble text lives inside the chapter. The retained file sources/stacks-more-algebra/retained-preamble.tex is a mechanical, line-for-line copy of the project's single shared preamble as staged in the snapshot (Xy-pic setup plus the native theorem declarations and mathematical macros). Nothing was added, removed or redefined, and the excerpt serves only as preamble, never as a span coordinate.\par\medskip
\noindent\textit{Editorial note.} The source URL is the source-level Stacks Project URL, as required when more than one card is produced from one source; the per-claim tag is carried by claim\_key and by this note. The counted statement and proof refer to no labelled item of the chapter and cite no bibliography entry, so this card contributes no external reference stub and needs no bibliography.\par\medskip
\noindent\textit{Allowed standard prerequisites (editorial).}\begin{itemize}
\item In an extension of discrete valuation rings every uniformizer of the smaller ring is a unit multiple of a power of a uniformizer of the larger ring, and the exponent is the ramification index. used at proof line 32973, where pi\_A = u pi\_B\textasciicircum{}e is written with u in the units of B. Provides the first of the two valuations to be compared.
\item If L over K is purely inseparable of characteristic p, then a power of every element of L lies in K, in particular a p-power power of pi\_B lies in K. used at proof lines 32974-32975, where pi\_B\textasciicircum{}q is placed in K for some p-power q and then written as v pi\_A\textasciicircum{}k. Supplies the second valuation equation whose exponent k is determined by the valuation of pi\_B\textasciicircum{}q, namely k = q/e.
\item Valuations of an extension of discrete valuation rings are multiplicative and take the value zero exactly on units, so equalities of elements up to units give equalities of valuations. used at proof lines 32976-32977, where taking valuations in B of pi\_A\textasciicircum{}q = u\textasciicircum{}q v\textasciicircum{}e pi\_A\textasciicircum{}\{ke\} gives ke = q. Converts the two writings of the same element into the numerical relation ke = q, hence e divides a power of p.
\end{itemize}\medskip
\begin{lemma}
\label{lemma-ramification-index-a-power-of-p}
Let $A \subset B$ be an extension of discrete valuation rings
inducing the field extension $K \subset L$. If the characteristic
of $K$ is $p > 0$ and $L$ is purely inseparable over $K$, then
the ramification index $e$ is a power of $p$.
\end{lemma}

\begin{proof}
Write $\pi_A = u \pi_B^e$ for some $u \in B^*$. On the other hand, we have
$\pi_B^q \in K$ for some $p$-power $q$. Write
$\pi_B^q = v \pi_A^k$ for some $v \in A^*$ and $k \in \mathbf{Z}$.
Then $\pi_A^q = u^q \pi_B^{qe} = u^q v^e \pi_A^{ke}$.
Taking valuations in $B$ we conclude that $ke = q$.
\end{proof}

\end{document}
